\chapter{Background and Related Work}\label{ch:background}

\needswork{STATUS: DRAFTED (October 2026) from the sources that are held and have been read.
Two gaps are marked where they occur and are gaps of \emph{reading}, not of writing: the
aggregate-signature schemes compared in \Cref{sec:bg-batch} are known here through one paper's
table, and no general FANET survey is held beyond the swarm-communication review cited below.}

\section{How citations are checked in this thesis}

No figure or number is quoted from a source that has not been read in full, and where a source
establishes prior art for something this thesis might otherwise have claimed as novel, that is
said explicitly. Every bibliography entry is compared by a script with its Crossref, DataCite or
arXiv record --- family names in order, initials, title and year --- and an entry with no
registry record must name the held document it was checked against.

Those rules exist because each was broken once. An early finding claimed this work was the
\emph{more optimistic} of two LoRa capacity models; it was drawn from the wrong figure of a paper
that reported pure-ALOHA results beside LoRa results, and was retracted: \textbf{quoting the
paper is not enough --- quote the figure}. And until October 2026 this section said that every
reference had been confirmed against its Crossref record, while four entries carried wrong
metadata and one paper was cited under the authors of a different one. The confirmation had
checked that each DOI resolved; it had not compared the author list. \textbf{A claim that
citations were checked belongs in a check that runs}, which is why the sentence is gone and the
script exists (\Cref{ch:reproducibility}).

\section{Flying ad-hoc networks and UAV telemetry}

A flying ad-hoc network is a set of unmanned aircraft that exchange information directly, without
a base station between them. Chen, Tang and Lao's review of swarm communication distinguishes
the architectures by where information is exchanged --- between the aircraft, or between the
aircraft and a central control centre --- and treats the routing protocols that each
needs~\cite{chen2020uavswarm}. This thesis is concerned with the simplest exchange the
decentralised architectures share, and the one all the others rest on: every aircraft
periodically \emph{broadcasts its own state} to whichever neighbours can hear it. There is no
route, no acknowledgement and no retransmission; a neighbourhood is a set of receivers within
range of one another, and \Cref{ch:sysmodel} models it as one collision domain.

\subsection{The workload, read at source}

What such a broadcast carries, and how often, is taken here from autopilot source code and from
a standard, not from secondary summaries.

\paragraph{PX4.} The MAVLink module configures its streams by link mode. For a link to a
companion computer (\texttt{MAVLINK\_MODE\_ONBOARD}) the global position is streamed at
$50$\,Hz, attitude at $100$\,Hz, the local position at $30$\,Hz and battery status at
$0.5$\,Hz~\cite{px4_mavlink_streams}; these four rates were read again in the source of release
v1.17.0 when it was built for \Cref{sec:e6}. The $50$ records per second of the adopted
operating point is that position rate.

\paragraph{ArduPilot.} Its default stream rates are set per vehicle type, and for the
multicopter firmware they default to zero~\cite{ardupilot_stream_rates}. An earlier draft of this work said ArduPilot ``defaults to
$1$\,Hz''; reading the parameter table corrected it. No single telemetry rate describes
``an autopilot'', which is why the operating point is treated as a variable in
\Cref{ch:codesign} and not as a constant.

\paragraph{The standard.} 3GPP specifies direct UAV-to-UAV local broadcast as a service, with
requirements that bound the same quantities this thesis optimises~\cite{3gpp_ts22125}: at least
ten messages per second, an end-to-end latency of at most $100$\,ms, and a payload of
$50$--$1500$\,B ``not including security-related message component(s)''. The last clause matters
here. The standard itself separates the bytes of authentication from the payload it sizes, and
says nothing about how large they may be; that silence is the gap this thesis measures.

\paragraph{How small the messages are.} Rotta and Mykytyn, designing authenticated telemetry
broadcast for a swarm, observe that MAVLink telemetry messages are ``much smaller than $256$
bytes''~\cite{rotta2024telemetrybroadcast}. The record of \Cref{ch:sysmodel} --- time, position,
velocity, battery, mode --- is \keyLean\,B standing alone and \deltaLean\,B as a difference from
its predecessor.

\subsection{How such networks are usually evaluated}

Simulation in ns-3 is the common instrument, and a recent swarm protocol stack shows what a
typical scenario states: $20$--$100$ nodes in a $400 \times 400 \times 1000$\,m volume, Random
Waypoint mobility, Nakagami fading and an 802.11ah MAC~\cite{bhatt2025uavswarmmac}. It does not
state the speed range or the pause time of its mobility model, the two parameters that define
what that model did. The observation is not a criticism of one paper. Cavalcanti \emph{et al.},
in an editorial note on a decade of vehicular-network research, report of an earlier survey of $151$
simulation papers that fewer than $15\%$ were repeatable and $87.5\%$ gave no confidence
intervals~\cite{cavalcanti2018vanet}. \Cref{ch:reproducibility} is this thesis's answer to that
record: every result is regenerated from stored runs, every simulated capacity carries an
interval, and the parameters that define each scenario are in the result files' headers.

Two things follow for the design of this study. Capacity is reported for \emph{one collision
domain with static nodes}, stated as such, because the contention model has no position term and
a mobility model would add parameters without adding a mechanism (\Cref{ch:validation}). And
the telemetry itself is synthetic, checked against public flight logs and against the autopilot
run in simulation (\Cref{sec:e6}), because a public corpus of $50$\,Hz swarm telemetry was not
found.

\section{Tamper-evident telemetry and hash-chained ledgers}

Linking each record to the hash of its predecessor makes an append-only log tamper-evident:
altering a record invalidates every subsequent link. The construction is Haber and
Stornetta's~\cite{haber1991timestamp}; Bitcoin made it the spine of a
blockchain~\cite{nakamoto2008}. UAV-blockchain systems, surveyed
in~\cite{mehta2020blockchainuav}, place telemetry or attestations on such a chain, gaining
integrity at a per-record overhead cost. This thesis targets exactly that substrate cost and is
agnostic to the consensus layer above it; finality and witnessing semantics are out of scope.

\subsection{What the chain gives and what the signature gives}

The two mechanisms answer different questions, and the threat model of \Cref{ch:sysmodel}
needs both.

A \emph{hash chain} fixes order and content relative to one value. Whoever holds the hash of the
latest record can detect any change to, removal of, or insertion among the records before it;
they need trust nothing but that one hash. The chain says nothing about \emph{who} wrote the
records: anyone can build a consistent chain over data of their choosing.

A \emph{signature} fixes origin. A record signed by a vehicle's key was produced by a holder of
that key, and a third party can be shown so. A signature on a single record says nothing about
the records around it: a signer can sign two different records for the same position in the
sequence, and each verifies.

Together they give what neither gives alone. A signed record that names its predecessor's hash
commits its signer to \emph{one history}; two validly signed records that claim the same
sequence number are then evidence of equivocation that anyone can check, because both carry the
signer's own signature. That property --- transferable evidence of inconsistency --- is why this
thesis requires a public-key signature and does not accept a shared-key tag in its place, and it
is the property lost by the cheaper per-message schemes of \Cref{sec:bg-permsg}.

It also explains a choice in the lean frame (\Cref{ch:sysmodel}). Within one frame the records
are signed together, so their order and content already rest on the signature; the link each
record would carry to its predecessor \emph{inside the same frame} can be recomputed by the
receiver and need not be sent. Only the link out of the frame, to the last record of the frame
before, carries information the signature does not. The stored ledger is unchanged.

\section{Signatures for constrained authentication}

ECDSA~\cite{fips186} yields $64$\,B P-256 signatures at 128-bit security but is
randomness-sensitive. Ed25519~\cite{ed25519,rfc8032} uses a twisted-Edwards curve with
deterministic nonces for fast, misuse-resistant, batchable verification. BLS~\cite{bls2001} gives
the shortest signatures via pairings, and~\cite{bgls2003} aggregates $n$ signatures into one at the
cost of expensive pairing verification.

AUTHBC introduces \emph{no new primitive}. It quantifies \emph{where and when} each scheme wins
once placement and batching are fixed (\Cref{ch:codesign}), finding Ed25519 dominant for
self-batching and BLS worthwhile only for cross-signer relay traffic.

Independent work on UAV swarms adopts precisely the design this thesis treats as its naive
baseline: Rotta and Mykytyn attach a HMAC-256 signature with a timestamp to \emph{each} broadcast
telemetry message, and observe that MAVLink telemetry messages are much smaller than
$256$\,bytes~\cite{rotta2024telemetrybroadcast} --- corroborating both the premise and the baseline
from outside this work. Their construction is a symmetric MAC under a group key, so it provides
integrity and group authenticity but \emph{not} non-repudiation, and therefore cannot underwrite a
ledger attributable to a specific UAV. That is a different trust model, not a cheaper solution to
this one.

\section{Authenticating messages one at a time, and signing streams}

\subsection{Per-message practice}\label{sec:bg-permsg}

MAVLink~2 signing appends a link identifier, a 48-bit timestamp and a 48-bit truncated SHA-256
tag computed with a shared secret~\cite{mavlink2signing}: thirteen bytes per message. It is cheap
and it authenticates a \emph{link} --- any holder of the key can produce any message --- so a
record cannot be attributed to one vehicle afterwards. Vehicular broadcast uses public-key
signatures: an ECDSA signature on each message, with the signer's certificate in some messages
and an eight-octet digest of it in the others~\cite{etsi_ts103097,ndss2024pqv2v}. That is the
trust model adopted here, and the per-message cost this thesis starts from.

\subsection{Amortising one signature over many packets}

This is a classical problem and the closest prior art to batching. Gennaro and Rohatgi sign the
first packet of a stream and chain the rest by hashes~\cite{gennaro1997streams}; one lost packet
breaks the chain. Wong and Lam sign a hash tree over a block and send in every packet the hashes
needed to reach the root~\cite{wong1999flows}, so each packet verifies alone --- the \emph{signing
cost} is amortised, the bytes are not. EMSS tolerates loss by carrying several hashes per packet
with a periodic signature packet, and TESLA replaces the signature by a MAC whose key is disclosed
later~\cite{perrig2000emss,perrig2002tesla}; TBRD applies TESLA to UAS Remote ID and reports
about half the overhead of signatures~\cite{veara2025tbrd}.

All of these spread \emph{one authenticator across packets}, and so make packets depend on each
other or on a later one. This thesis does the opposite: it puts \emph{several records in one
packet}, which is possible only because a telemetry record is a small fraction of a frame. Each
frame then stays self-verifying and the cost moves to latency, at the sender. TESLA instead delays
verification at the receiver, and a disclosed MAC key gives no non-repudiation, which a ledger
needs. The lesson of the stream-signing literature that does transfer is the one this thesis
learned late: a unit that depends on its predecessor inherits its predecessor's loss
(\Cref{thm:t3prime}). \Cref{sec:stream} places each of these schemes at this thesis's operating
point, by authenticator bytes and by frames per second.

\subsection{Hash-linked records}

Linking each record to the hash of its predecessor to make a sequence tamper-evident is due to
Haber and Stornetta~\cite{haber1991timestamp} and is the chain in every blockchain. The LoRaWAN
link layer whose payload limits bound \Cref{ch:lowrate} is specified
in~\cite{lorawan_ts001,lora_rp002}.

\subsection{Certificates}

A signature is checked against a public key, and a receiver must learn that key from somewhere.
Vehicular practice attaches the signer's certificate to some messages and a short digest of it
to the rest~\cite{etsi_ts103097,ndss2024pqv2v}; \Cref{ch:codesign} charges that policy to every
frame. IEEE~1609.2 defines two kinds of certificate. An \emph{explicit} certificate carries the
verification key and the issuer's signature over it. An \emph{implicit} certificate (ECQV)
carries only a reconstruction value from which the key is derived using the issuer's root
certificate, and is shorter: a signed message is at most $226$\,B with one against $330$\,B with
an explicit certificate~\cite{ndss2024pqv2v}. This thesis prices explicit certificates only. The
certificate figures of \Cref{ch:codesign} are therefore upper bounds, for the baseline and the
design alike.

\section{Batch and aggregate verification}\label{sec:bg-batch}

Zhang \emph{et al.}~\cite{zhang2008batch} let a roadside unit verify multiple received signatures
simultaneously, motivated by the same pressure from the opposite side: a hard deadline --- DSRC's
$300$\,ms --- within which sequential verification does not scale. Their batching is
\emph{receiver-side} and buys verification time; the batching in this thesis is \emph{sender-side}
and buys airtime. The two are complementary rather than competing. The nearest FANET work is
of the same kind: Rajasekaran \emph{et al.}~\cite{rajasekaran2022batchfanet} authenticate a
group of drones to a user in one batch check of a pairing-based scheme, and report the communication cost of
$n$ authentications as $n$ times that of one ($1184\,n$ bits) --- the verifier's work is
shared, the bytes are not. This work generalises from
receiver batching to authentication \emph{placement}, co-designs it with encoding and scheme, and
adds a loss-robustness frontier that pure aggregation does not address.

Certificateless aggregate signature (CLAS) schemes for vehicular broadcast are the closest
external comparator. Li \emph{et al.}~\cite{li2025clas} tabulate the per-message size of their
own scheme and four others, for a $67$\,B traffic message: $583$\,B (theirs, and Cahyadi
\emph{et al.}), $596$\,B (Xu \emph{et al.}), $735$\,B (Liang \emph{et al.}) and $859$\,B (Wang
\emph{et al.}) --- the last four as reported there, not read at source. In every one the cost of
$n$ messages is $n$ times the cost of one. \textbf{The finding is the axis, not a ratio}:
aggregation compresses the verifier's work, while each message still carries its own identity,
key and signature components; batching compresses what is sent. No claim is made to beat those
schemes. They provide conditional privacy, which this work does not, their group elements are
$128$\,B against a $64$\,B signature here, and their figures are for a different message.

What these schemes are for is worth stating, since it is not what this thesis is for. In a
certificateless scheme a key-generation centre issues part of each signer's private key and the
signer adds a secret of its own, which removes certificates without giving the centre the whole
key; the aggregate property lets a verifier check many signatures at once. Li \emph{et al.}
show that one such scheme, by Cahyadi \emph{et al.}, falls to a public-key replacement attack
and to a malicious key-generation centre, and repair it~\cite{li2025clas}. That is the
literature's centre of gravity: constructions, attacks on them and proofs, with communication
cost reported as a per-message size and never as what a channel can then carry.

\needswork{A gap of reading, stated: the four schemes compared in~\cite{li2025clas} are known
here only through that paper's table and were not read at source. A fuller treatment of their
security models needs those papers.}

\section{Compact serialisation and delta coding}

CBOR~\cite{rfc8949} and MessagePack~\cite{msgpack} give compact binary encodings. Delta coding of
telemetry time series is established practice, most prominently in Gorilla~\cite{pelkonen2015gorilla},
which stores delta-of-delta timestamps and XOR-folded values for roughly a $10\times$ reduction.

The variant used here differs in one respect the link forces. Gorilla is a database and can chain
deltas indefinitely. On a lossy broadcast link a record coded against one the receiver never got
cannot be decoded, and a frame that cannot be decoded cannot be verified (\Cref{thm:t3prime}).
The first design of this work re-anchored on a self-contained record every $K = 16$ records, and
missed its own verifiability target for that reason; the design reported here starts
\emph{every frame} with a self-contained record and codes only the records behind it in the same
frame (\Cref{sec:e7}).

\section{Channel models for IEEE 802.11}

Bianchi's fixed-point analysis of the DCF~\cite{bianchi2000} is the standard model for saturated
unicast throughput. Broadcast is \emph{not} that model with the acknowledgement removed --- a point
Ma and Chen make explicitly, and which this project rediscovered expensively. Their broadcast
model~\cite{machen2007,machen2008} accounts for the backoff-counter Consecutive Freeze Process,
which dominates when the contention window is small relative to the station count. Tinnirello,
Bianchi and Xiao~\cite{tinnirello2010} supply the anomalous-slot refinement that explains a
residual small-frame bias observed in \Cref{ch:validation}.

\section{LoRa and LoRaWAN capacity}

Haxhibeqiri \emph{et al.}~\cite{haxhibeqiri2017lora} provide the only published LoRa capacity model found that is
grounded in hardware interference measurements rather than an analytical collision assumption, and
state it in closed form so it can be evaluated at this work's operating point rather than merely
quoted at theirs. Zirak \emph{et al.}~\cite{zirak2021swarm} supply the only air-to-air LoRa PDR-versus-range
field data located, which range-limits the capacity result of \Cref{ch:lowrate}.

\subsection{Limits that are rules, not physics}

Two of the limits that shape \Cref{ch:lowrate} are set by documents and not by the radio. The
LoRa physical layer carries up to $255$\,B in a frame~\cite{semtech_sx1276}; the payload limits
of $51$, $115$, $222$ or $242$\,B that the exclusion bound is computed against are those of
LoRaWAN's regional parameters for the EU863--870 band, one per data rate~\cite{lora_rp002}. A
system that used the LoRa radio without the LoRaWAN link layer would not be bound by them, and
the exclusion is stated with that condition. The second is the duty cycle: the same document
limits a device's transmit time in the band to less than $1\%$, which fixes how many records per
second a node may send whatever its neighbourhood. The link layer itself is specified
in~\cite{lorawan_ts001}, which also requires an end device to vary its transmit times
pseudo-randomly ``to prevent systematic synchronization of populations of end-device
transmissions'' --- the standard's own statement of the artifact that strictly periodic
simulated senders produce (\Cref{ch:reproducibility}).

\section{Where this thesis sits}

The adjacent literature almost uniformly proposes a \emph{new construction} and evaluates it
against its predecessors. This thesis deliberately proposes none. Its contribution is a boundary:
the links on which a frame that verifies alone cannot be sent at any batch size, stated with the
conditions under which that holds. For a given authenticator the result is arithmetic rather
than empirical, so unlike a performance figure it does not move as models or hardware improve. It
does move with the authenticator --- a $48$\,B signature or a $13$\,B tag fits where a $64$\,B
signature does not --- and \Cref{ch:lowrate} gives the matrix. It is invisible to the byte
models the literature optimises against.

\needswork{The bibliography holds every source that is held and has been \emph{read}, and stops
there deliberately. Reaching a higher count requires a genuine sourcing pass --- find, download,
read, then cite. Padding the list would be the same defect this project spent its audit
removing. The reading most worth doing next: a general FANET survey; the four aggregate-signature
schemes named above; and the stream-signing schemes' successors.}
